% !TEX root = main.tex

\section{Explicación respuestas}

\begin{enumerate}

\item Que ocurre en un protocolo por turnos si un nodo retiene el token mas tiempo del permitido?
Sin mecanismo de protección: los demás nodos sufren hambruna(starvation), con THT (Token holding timer), el token se libera al vencer el timer aunque el nodo no haya terminado.

\item Que problema mitiga RTS/CTS en WiFi?
El problema del nodo oculto: dos estaciones que no se escuchan entre sí pero si al AP \textbf{Access point} pueden transmitir simultáneamente causando colisión; RTS/CTS reserva el canal alertando a todos los nodos en rango del AP.

\item Por qué la dirección MAC no cambia cuando el host cambia de red?
Porque la MAC está grabada en la ROM de la NIC (hardware) e identifica la interfaz, no la red; lo que cambia es la IP asignada por el nuevo DHCP.

\item En que tipo de red es preferible un esquema basado en turnos frente a CDMA?
En redes \textbf{industriales} con nodos fijos que requieren latencia máxima garantizada (tiempo real duro) como PLCs en una linea de ensamblaje.

\item Por qué CDMA puede considerarse mas flexible que un esquema por turnos?
CDMA permite acceso simultáneo sin coordinación centralizada: múltiples usuarios transmiten en el mismo canal usando códigos ortogonales; agregar o quitar nodos no requiere reconfigurar turnos.

\item En qué aplicación real es el CRC indispensable?
En discos duros y SSDs; ya que cada sector incluye in CRC para detectar corrupción silenciosa al leer.

\item Que mejoras trajo 802.11ac respecto a 802.11n en terminos de agregación de canales?

802.11ac amplió los cnales de 40 MHz hasta 80 y 160 MHz, solo en 5 GHz, añadió 256-QAM, 8 streams espaciales y MU-MIMO downlink, velocidad teorica hasta 6.93 Gbps.

\item En que consiste un ataque de ARP Spoofing?
El atacante envía ARP replies falsos no solicitados mapeando su MAC a la IP de la victima o del gateway, convirtiendose en man-in-the-middle.

\item Por qué el checksum puede fallar al detectar ciertos patrones de error?

Porque si dos errores se cancelan en la suma (uno incrementa en k y otro decrementa en k), el total no cambia

\item Cuales son dos servicios que ofrece la capa de enlace en una red local?
Entramado (framing) y control de acceso al medio (MAC)

\item Por qué ARP no requiere un servidor centralizado?
Porque usa el broadcast de la LAN local; la pregunta llega a todos los nodos del segmento y el dueño de la IP responde directamente, \textbf{sin intermediario}

\item Que caracteristica hace a 802.11be (WIFI 7) especialmente relevante para aplicaciones multimedia?

Multi-link operation (MLO) transmite y recibe simultaneamente en multiples bandas (2.4, 5 y 6 GHz) garantizando latencia ultra-baja y alta resiliencia para streaming 8K, VR/AR y gaming en la nube.

\item Como mejora RTS/CTS el comportamiento de CSMA/CA cuando existe un nodo oculto?

El emisor envía un RTS pequeño al AP, el AP responde con CTS en broadcast; todos los nodos que escuchan el CTS (incluido el nodo oculto) saben que el canal está reservado y esperan.

\item Por qué 802.11 usa CSMA/CA en lugar de CSMA/CD?

Porque una radio no puede transmitir y recibir simultaneamente en la misma frecuencia, no puede detectar colisiones mientras transmite, por eso previene CA en lugar de detectar CD.

\item Por qué se requiere trunking para interconectar varios switches con multiples VLANs?

Para transportar tramas de varias VLANs por un solo enlace fisico etiquetandolas con 802.1Q, evitando necesitar un cable fisico separado por cada VLAN.

\item Como mejora una VLAN la seguridad en una red corporativa?

Separando el tráfico en dominios de broadcast aislados; las tramas de finanzas no llegan a RRHH sin pasar por el router (que aplica ACLs/firewall)

\item que hace un switch cuando recibe una trama con destino desconocido?
La inunda(flood) por todos los puertos excepto el de entrada, garantizando que le destino la reciba y permitiendo al switch aprender el puerto del destino cuando responda.

\item por qué los switches pueden operar en modo full-duplex mientras los hubs no?
porque cada enlace switch-host es punto a punto dedicado sin bus compartido; el switch puede recibir de A mientras transmite a B simultaneamente, eliminando colisiones y duplicando el thorughput efectivo del enlace.

\item En CSMA/CD, que sucede cuando dos hosts detectan simultaneamente una colisión?

Ambos paran, emiten una señal jam de 32 bits y aplican el algoritmo de backoff exponencial binario.

\item cuales son dos diferencias clave entre 802.11b y 802.11g en velocidad y compatibilidad?

802.11b: 11 Mbps con DSSS en 2.4 GHz. 802.11g: 54 Mbps con OFDM en 2.4 GHz y retrocompatible con 802.11b (aunque la presencia de clientes b reduce el throughput de los clientes g por el mecanismo de protección)

\item Por qué CSMA/CD funciona en ethernet tradicional pero no en redes wi-fi?
en ethernet, la NIC puede escuchar el cable mientras transmite y detectar colisiones. En wi-fi la radio no puede transmitir y recibir simultaneamente en la misma frecuencia.

\item Por qué CRC es adecuado para enlaces ruidosos como redes inalambricas?

Porque usa división de polinomios en GF(2) y detecta garantizadamente todas las rafagas de error de hasta r bits.

\item por qué la trama 802.11 contiene mas direcciones que la trama ethernet?

porque involucra 3 actores (cliente, AP, destino en red cableada): Addr1=AP receptor, Addr2=cliente transmisor, Addr3=destino final en la red cableada; el AP necesita Addr3 para saber a quién reenviar tras desencapsular

\item cual es el rol de la NIC en el funcionamiento de la capa de enlace?
encapsula datagramas en tramas, calcula el FCS, gestiona CSMA y almacena la dirección MAC en ROM.

\item cual es la función del preámbulo de ethernet?
permite a la NIC receptora sincronizar su reloj con el transmisor (7 bytes) y marca el inicio de la trama como el SFD (1 byte, 10101011)

\item que funcion cumple el tag 802.1q en una trama ethernet?
identifica a que VLAN pertenece la trama mediante el VID de 12 bits

\item que información transporta el beacon frame de un AP?
B
SSID, BSSID (MAC del AP), timestamp de sincronización, capacidades soportadas (velocidades, seguridad WPA2/WPA3), canal activo, intervalo de beacon y TIM (mapa de tráfico para modo ahorro de energía)

\item sistema con paridad par. Transmisor envía 1101001. Receptor recibe 1100001. Se detecta el error?
Si, el dato recibido tiene 3 unos (impar) y la paridad par esperaba cantidad par

\item En 802.11ax (Wi-Fi 6), ¿cómo ayuda OFDMA a mejorar la eficiencia en redes densas?
B
Dividiendo el canal en Resource Units (RUs) asignables a múltiples clientes simultáneamente en la misma transmisión, en lugar de que un solo cliente use todo el canal como en OFDM

\item en qué escenario el FCS detecta un error que ni paridad ni checksum detectarían?
Una ráfaga de interferencia electromagnética que altera 20 bits consecutivos: paridad (número par, no detecta) y checksum (pueden cancelarse) fallan, pero CRC-32 detecta garantizadamente todas las ráfagas <= 32 bits

\item qué es un BSS y cómo se relaciona con un AP?
Basic Service Set: unidad fundamental de 802.11 formada por un AP (nodo central, BSSID = su MAC) y todas las estaciones asociadas a él; toda comunicación entre clientes pasa por el AP

\item que significa que un checksum sea de complemento a uno?
Que el checksum es la inversión de todos los bits de la suma, y la verificación da todo unos si no hay error

\item por qué la falta de control de congestión en Ethernet no es un problema en redes conmutadas modernas?
porque TCP detecta las pérdidas (timeout o ACKs duplicados) y reduce su ventana de congestión, gestionando la congestión en la capa de transporte

\item por qué 802.11n introdujo MIMO y cómo impactó en el rendimiento de Wi-Fi?
MIMO usa múltiples antenas para enviar varias corrientes de datos espaciales simultáneamente (spatial multiplexing) y para diversidad anti-multipath; pasó de 54 Mbps (802.11g) a hasta 600 Mbps teóricos con 4 streams en 40 MHz

\item por qué un enlace punto a punto como PPP no necesita mecanismos de arbitraje del medio?
porque solo hay dos nodos en el enlace, por lo que no puede haber colisión

\item por qué el checksum es más robusto que la paridad simple?
Porque opera palabra a palabra (grupos de 16 bits) y detecta errores que afectan el valor numérico, incluyendo patrones de error doble que la paridad ignora

\item compare token passing con CDMA en términos de eficiencia en redes con baja carga.

CDMA es más eficiente con baja carga: cualquier nodo transmite al instante. Token desperdicia tiempo haciendo circular el token por nodos inactivos
\item nn switch recibe una trama en un puerto cuyo origen ya está registrado en otro puerto. ¿Qué ocurre?
El switch actualiza la tabla MAC con el nuevo puerto, asumiendo que el host se reconectó o cambió de ubicación

\item en qué escenario token passing reduce la latencia respecto a CDMA?
En redes con alta carga y transmisión de bloques grandes: token garantiza el canal completo sin interferencia MAI, mientras CDMA con muchos usuarios simultáneos degrada por interferencia de acceso múltiple

\item como usa un host A el protocolo ARP para conocer la MAC de un host B?
A envía un ARP Request en broadcast (FF-FF-FF-FF-FF-FF) con la IP de B; B responde con un ARP Reply unicast con su MAC; A guarda la entrada en su caché

\item qué significa que un error sea 'divisible por el polinomio generador'?
Que el patrón de error E(x) es múltiplo de G(x), por lo que el residuo de la división sigue siendo 0 y el error no se detecta

\item por qué Ethernet se considera una tecnología de 'mejor esfuerzo'?
Porque no garantiza entrega, orden ni latencia: si el switch descarta una trama por buffer lleno, no hay retransmisión a nivel de enlace

\item qué sucede si el campo de tipo/longitud de la trama Ethernet no concuerda con el tamaño real del payload?
el receptor puede leer bytes de la trama siguiente (runt frame) o ignorar datos útiles; el FCS normalmente detecta la inconsistencia y descarta la trama

\item por qué los switches permiten dominios de colisión más pequeños que los hubs?
Porque el switch conmuta trama a trama y cada puerto es su propio dominio de colisión dedicado (full-duplex posible), mientras el hub crea un único dominio de colisión para todos los puertos

\item En CSMA/CA, ¿cómo ayuda el mecanismo de backoff aleatorio a reducir colisiones?
Cada nodo sortea un contador aleatorio y lo decrementa solo mientras el canal está libre; el primero en llegar a 0 transmite, los demás congelan su contador

\item qué sucede si una caché ARP contiene una entrada obsoleta?

Las tramas se envían a la MAC antigua: el host incorrecto las recibe o se pierden hasta que la entrada expire y se refresque con un nuevo ARP

\item por qué ARP se implementa en la frontera entre la capa de enlace y la capa de red?
Porque traduce IPs (capa 3) a MACs (capa 2): la capa de red necesita la MAC para construir la trama

\item en qué escenario el uso de paridad simple es suficiente?
En teclados PS/2 o interfaces UART de baja velocidad en entornos controlados donde la probabilidad de error doble es despreciable

\item ¿En qué escenario la detección de errores es fundamental aunque la corrección se haga en capas superiores?

En streaming de video UDP sobre Wi-Fi: detectar la trama dañada evita entregar basura a la aplicación

\item ¿Cómo aprende un switch la tabla MAC cuando está inicialmente vacía?
Al recibir una trama, aprende la MAC origen y su puerto de entrada; si la MAC destino es desconocida, inunda por todos los demás puertos; las respuestas completan la tabla

\item ¿Por qué se afirma que la capa de enlace está 'cercana al hardware'?
Porque su implementación reside en la NIC (hardware/firmware), no en el procesador principal

\item ¿Qué problema se evitaría al usar VLANs en una red universitaria con miles de usuarios?
Tormentas de broadcast: sin VLANs, un ARP Request se replica a los miles de hosts del segmento; con VLANs, el broadcast queda confinado a grupos pequeños

\item ¿Qué ocurre si un switch recibe una trama tagged 802.1Q en un puerto configurado como access?
La descarta para prevenir ataques de VLAN hopping (double-tagging)

\item ¿Por qué Wi-Fi usa CSMA/CA en lugar de CSMA/CD? Describe el mecanismo.
Porque no puede detectar colisiones; en cambio previene: escucha el canal, espera DIFS, aplica backoff aleatorio, transmite y espera ACK

\item ¿Cuándo la paridad simple NO puede detectar un error?
Cuando se alteran dos bits (número par de errores)
















\end{enumerate}